\section{Abordagem Proposta}
\label{sec:metodo}

A abordagem tem dois componentes, construídos sobre a mesma segmentação da transcrição em turnos de
fala (Seção~\ref{sec:dados-particao}). O primeiro, a Unidade Deliberativa Verificável (UDV), liga cada
uma das 2.203 opiniões publicadas a um trecho da fala da própria pessoa na transcrição, com a posição
exata, o tipo de ligação encontrada e a origem dessa ligação. O segundo reúne tudo o que cada pessoa
disse nas audiências de que participou e gera, com um LLM, um perfil textual escrito apenas a partir
dessas falas.

\subsection{Vis\~ao Geral da Arquitetura}
\label{sec:metodo-visao}

\input{figures/architecture}

A Figura~\ref{fig:arquitetura} mostra o fluxo: cada transcrição do LDS é dividida em turnos, o
cabeçalho de cada turno define o autor do trecho, e a partir daí o processamento segue em dois ramos. No ramo da UDV
(Seção~\ref{sec:metodo-udv}), cada audiência é tratada isoladamente: os participantes citados na
matéria são associados aos seus turnos, e cada opinião recebe a sentença da fala da pessoa encontrada
por citação literal ou por similaridade semântica, com offsets de caracteres e um nível de confiança
declarado. No ramo dos atores (Seções~\ref{sec:metodo-atores} e~\ref{sec:metodo-perfis}), os turnos
são agrupados por pessoa entre audiências, filtrados, restritos às audiências de treino e entregues a
um LLM que escreve o perfil. Os dois ramos resolvem identidade de formas diferentes: a UDV associa o
nome da matéria aos cabeçalhos de uma audiência, e o ramo dos atores associa cabeçalhos de audiências
diferentes entre si. Eles compartilham as chaves (audiência, índice do turno e offsets na transcrição),
de modo que a evidência de uma UDV cai dentro de um turno de um registro de ator e os dois artefatos
podem ser cruzados.

\subsection{Unidades Deliberativas Verificáveis}
\label{sec:metodo-udv}

As opiniões do PublicHearingBR registram o que a matéria atribui a cada pessoa, e nada no conjunto
indica de que ponto da transcrição elas saíram (Seção~\ref{sec:dados}). Sem esse vínculo, não há como
mostrar ao leitor a passagem que originou uma opinião, conferir se a fala a sustenta nem estudar o que
uma pessoa defendeu a partir do que ela disse. A UDV é um registro JSON por opinião que cria esse
vínculo (Quadro~\ref{lst:udv}). No exemplo, a matéria escreve ``presidente da comissão, que é
eleita'' e a transcrição, ``presidente que foi eleita''; o prefixo da citação casa literalmente, e a
diferença posterior fica visível na comparação entre a opinião e o texto da evidência.

\textbf{Resolução da pessoa.} O nome no LDS é o da matéria, que costuma diferir do cabeçalho
(``Rodrigo Marinho'' na matéria, ``RODRIGO SARAIVA MARINHO'' na transcrição), e quem preside a sessão
aparece como \textit{PRESIDENTE}, com o nome entre parênteses. Cada turno gera candidatos de nome (o
nome do cabeçalho e, na presidência, o nome entre parênteses), e um participante casa com um candidato
quando os conjuntos de palavras, sem acento e em maiúsculas, são iguais ou um contém o outro, desde que
ambos tenham duas palavras ou mais. Um nome de uma palavra só casa com um candidato idêntico ou, na
falta dele, com o único falante da audiência que tenha aquela palavra no nome. Com essa regra, 1.020
dos 1.065 participantes recebem ao menos um turno; os 45 restantes somam 90 opiniões.

\begin{lstlisting}[float=tbp,caption={Uma UDV. Os offsets de caracteres apontam para a transcrição da
audiência 1, e o turno de fala é identificado pelo seu índice na transcrição. Os campos que registram o
codificador e o corte usados foram omitidos.},label=lst:udv]
{
  "id": "udv-1-3-4",
  "hearing_id": 1,
  "actor": {"name": "Arlindo Chinaglia", "role": "Deputado (PT-SP)"},
  "proposition": "\"A resposta do Musk foi atacar a presidente da comissão, que é eleita por todos os países da Comunidade Europeia.\"",
  "evidence": {
    "text": "A resposta do Musk foi atacar a presidente que foi eleita por todos os países da comunidade europeia, inclusive a democrática Hungria.",
    "support_type": "direct_quote",
    "score": null,
    "quote_prefix": "A resposta do Musk foi atacar",
    "start_char": 58193, "end_char": 58327, "speaker_turn": 27
  },
  "tier": "quote_found",
  "provenance": "weak"
}
\end{lstlisting}

\textbf{Sentenças candidatas.} A fala da pessoa é dividida em sentenças dentro de cada turno, e
cada sentença guarda o índice do turno de origem. Uma sentença nunca atravessa a fronteira entre dois
turnos, onde colaria o fim de uma fala ao começo de outra, com falas de terceiros no meio. Ficam as
sentenças com quatro palavras ou mais que não sejam só marcação de palco, como ``(Palmas.)'': 115.599
sentenças nas 206 audiências.

\textbf{Citação literal.} Muitas opiniões trazem trechos entre aspas, que a matéria apresenta como
fala literal, e a citação permite uma ligação que não depende de modelo. De cada citação com 10
caracteres ou mais saem prefixos de 10, 6, 4 e 3 palavras, procurados nessa ordem dentro de cada turno
da pessoa, sem distinção de maiúsculas e sem casar no meio de uma palavra; entre as citações de uma
opinião vence a de prefixo mais longo. Um prefixo de seis palavras ou mais faz da sentença que o
contém a evidência, com o nível \texttt{quote\_found}; se o prefixo ocorre mais de uma vez na fala, vence
a sentença com maior índice de Jaccard de palavras com a opinião. Prefixos de três ou quatro palavras
costumam ser expressões comuns, que casam em pontos da fala sem relação com a citação. O limite de seis
palavras veio de uma leitura exploratória dos casamentos agrupados por tamanho de prefixo, sem
protocolo de anotação. Das 2.203 opiniões, 960 têm alguma citação, 549 têm algum prefixo encontrado na
fala da pessoa e 277 têm prefixo de seis palavras ou mais.

\textbf{Similaridade semântica.} Quando não há citação confiável, a opinião e as sentenças
candidatas da pessoa são codificadas pelo Serafim 335m~\citep{gomes2024serafim}, um codificador de
sentenças para o português (vetores de 1.024 dimensões, entrada limitada a 128 tokens), e a evidência
é a sentença de maior cosseno com a opinião. Um codificador denso foi preferido ao TF-IDF porque a
matéria costuma parafrasear a fala, e a comparação por vocabulário compartilhado falha nesses casos.
Se um prefixo curto (menos de seis palavras) de uma citação da opinião cai na sentença escolhida, o
tipo de suporte registra a coincidência, que é uma corroboração fraca e não muda o nível. Cada UDV
guarda a versão do codificador e o corte com que foi produzida.

\textbf{Nível, proveniência e posição.} O nível resume o tipo de ligação encontrada
(Tabela~\ref{tab:udv-niveis}), e os níveis \texttt{semantic\_match\_high} e
\texttt{semantic\_match\_weak} dizem só de que lado do corte ficou o cosseno da melhor sentença, sem
que nenhum deles afirme que a sentença sustenta a opinião. A proveniência registra a origem da ligação: \texttt{weak} (supervisão fraca,
rótulo produzido por regra, sem conferência humana) para o casamento de texto e \texttt{model} para o
codificador; nenhum registro tem proveniência humana. A sentença escolhida é localizada na transcrição
original por uma expressão regular restrita ao turno de origem, e as 2.105 evidências têm offsets
localizados. Os níveis \texttt{no\_evidence} e \texttt{person\_not\_resolved} registram que a fala não
foi encontrada e não indicam erro da matéria: os 8 registros \texttt{no\_evidence} são de participantes cuja única fala transcrita é a marcação
``(Manifestação em LIBRAS.)'', e entre as pessoas não resolvidas há grafias diferentes do nome,
instituições citadas como participantes e pessoas que não falaram.

\input{tables/udv-tiers}

\textbf{Calibração do corte.} O corte que separa \texttt{semantic\_match\_high} de
\texttt{semantic\_match\_weak} foi calibrado só nas audiências de treino, porque a amostra de validação
humana é sorteada no teste e depende dessa fronteira. A regra declarada antes de qualquer resultado
escolheria o corte que melhor separa acertos e erros do codificador em opiniões com os trechos entre
aspas removidos. Ela se mostrou degenerada: só 12 das 142 consultas do treino têm a sentença da
citação em primeiro lugar, e os erros têm cosseno mais alto que os acertos (medianas 0,655 e 0,578).
O corte adotado usa como positivos as 183 opiniões do treino com citação confiável, cada uma contra a
sentença da própria citação, e como negativos a mesma opinião contra uma sentença sorteada de outra
audiência do treino. O corte é o ponto médio entre o primeiro quartil dos positivos (0,628) e o
terceiro quartil dos negativos (0,281): 0,454, arredondado para 0,45 (IC 95\% por \textit{bootstrap}
de audiências, de 0,437 a 0,469). Com negativos tirados da fala da própria pessoa, que é a escolha que
o codificador faz de fato, a mesma regra daria 0,50. O corte define uma fronteira entre rótulos e não
foi validado como medida de confiança.

\textbf{Verificação.} Um verificador separado refaz, a partir do LDS e com as mesmas funções, a
resolução de pessoa, as sentenças candidatas e a busca de citação de cada opinião, e confere a
coerência de cada registro com as regras: opinião e ator iguais aos do LDS, texto da evidência dentro
de um único turno do ator, texto recuperado pelos offsets igual ao gravado, e nível coerente com o tipo
de suporte, o cosseno e o corte. Ele não recalcula a escolha do codificador nem julga se a evidência
sustenta a opinião. O conjunto de UDVs descrito aqui passa sem nenhuma inconsistência, e uma versão
anterior do conjunto, regenerada com o código que a produziu, saiu idêntica byte a byte à original.

\subsection{Falas por ator entre audiências}
\label{sec:metodo-atores}

Como a UDV trata uma audiência de cada vez, descrever o que uma pessoa defende ao longo do tempo exige
reunir tudo o que ela disse em todas as audiências de que participou e, para isso, decidir quando dois
cabeçalhos de audiências diferentes são a mesma pessoa. A regra de nome contido da UDV é
aceitável dentro de uma audiência, que tem poucos falantes, mas entre audiências junta pessoas
diferentes (``Luiz Lima'' e ``Gustavo Luiz de Lima Correia''). O registro de ator usa como chave o
nome do cabeçalho sem acento e em maiúsculas (na presidência, o nome entre parênteses), mais uma lista
de mesclas. Os 153 pares de nomes candidatos, em que um nome contém o outro ou os dois têm as mesmas
palavras, foram revisados um a um com evidência do próprio conjunto: o cargo na matéria, a
apresentação feita pelo presidente, a auto-apresentação, o partido e a UF do cabeçalho e a data da
audiência. A revisão resultou em 37 grupos de mescla, uma chave dividida por audiência (um homônimo
exato) e um par sem resolução, mantido separado. Nenhuma identidade foi conferida contra registros
externos ao conjunto.

Três regras retiram turnos que não trazem a fala da própria pessoa ou que só conduzem a sessão. Os
turnos de intérprete e de tradução simultânea saem, porque a fala pertence a outro participante. Os
turnos de presidência com menos de 50 palavras também saem. Quem preside fala muito: 42,7\% das
palavras dos falantes presentes em duas ou mais audiências estão em turnos de presidência, e boa parte
disso é passar a palavra, controlar o tempo e agradecer. Descartar a presidência inteira, porém,
perderia opinião real, porque 336 das 2.105 evidências de UDV estão em turnos de presidência. O turno
de presidência mediano tem 28 palavras; o corte de 50 mantém 2.264 dos 6.266 turnos de presidência,
com 89\% das palavras, e 332 das 336 evidências. Os turnos mantidos ficam marcados como de presidência.
Por fim, saem turnos vazios e turnos formados só por marcação de palco. Ficam 13.190 dos 17.264 turnos,
e o texto de cada um é copiado sem limpeza; a igualdade com o intervalo de offsets da transcrição foi
conferida nos 13.190. Os 301 atores que falam em duas ou mais audiências são o material dos perfis
(Quadro~\ref{lst:ator}).

\begin{lstlisting}[float=tbp,caption={Registro de falas por ator, restrito às audiências de treino. Os textos
foram truncados em {[\ldots]}, e o campo com a concatenação dos turnos de cada audiência foi omitido.},
label=lst:ator]
{
  "actor": "DANIEL BARBOSA",
  "has_party_header": true,
  "party_uf": ["Bloco/PP - AL"],
  "hearings": [
    {"hearing_id": 31, "turns": [
      {"turn_index": 4, "role": "speaker", "start_char": 23770, "end_char": 25429,
       "text": "Obrigado, Deputada Socorro Neri.\n\nAs palavras de V.Exa. foram muito bem colocadas. [...]"}]},
    {"hearing_id": 168, "turns": [
      {"turn_index": 17, "role": "speaker", "start_char": 73694, "end_char": 74807,
       "text": "Bom, Presidente Alice, eu vou fazer uma fala rápida. Eu fui a favor da retirada do Ministério da Ciência, Tecnologia e Inovação do teto de gastos [...]"}]}
  ]
}
\end{lstlisting}

\subsection{Perfis por ator}
\label{sec:metodo-perfis}

A fala de um ator se espalha por audiências de temas diferentes e se mistura com condução e
cortesia. O perfil condensa essa fala num texto que
descreve o que a pessoa defende, critica, propõe e cobra, escrito por um LLM apenas a partir das
falas. Cada perfil é gerado só com as audiências de treino (Seção~\ref{sec:dados-particao}), para que
nenhuma fala ou opinião das audiências de teste chegue ao prompt.

Cada ator recebe uma única geração, e o prompt do usuário traz todas as suas falas de treino em blocos
por audiência, em ordem cronológica; cada bloco começa com a data da
matéria (Seção~\ref{sec:dados-particao}) e o assunto registrado no LDS, e os turnos de presidência
vêm precedidos da marca \texttt{[presidência da sessão]}. A data permite situar mudanças de posição
no tempo, e o assunto ajuda a desfazer referências vagas como ``este projeto de lei''. O prompt de
sistema impõe três grupos de regras; os dois prompts estão no Apêndice~\ref{ap:prompt}.

\textbf{Fonte única.} O perfil usa somente as falas, sem conhecimento prévio sobre a pessoa,
partidos, organizações ou temas. O cabeçalho da audiência vem da matéria e serve só para situar as
falas; nada que apareça apenas nele é atribuído à pessoa. Siglas não são expandidas, nomes e cargos
de pessoas citadas não são completados, e nenhum fato recebe uma data que a fala não traz. Dados
citados entram como afirmação da pessoa (``afirma que''). Proposta ou opinião de terceiros que a
pessoa só relata, resume ou questiona não vira posição dela. Perguntas, hipóteses e condicionais não
viram afirmação, e números mantêm o recorte a que se referem. Falas escassas ou protocolares geram
um perfil de duas a quatro frases.

\textbf{Condução e cortesia.} Em qualquer turno, passar a palavra, controlar o tempo, citar
requerimentos, cumprimentar, agradecer e elogiar são desconsiderados; dos turnos de presidência se
aproveita só o que é posição própria.

\textbf{Escrita.} Prosa corrida em terceira pessoa, sem listas nem títulos, começando pelo
conteúdo mais distintivo da pessoa. A posição concreta (``propõe que as operadoras sejam obrigadas a
bloquear celulares roubados'') é preferida ao rótulo abstrato (``preocupado com segurança
pública''), com o vocabulário do domínio presente nas falas. Constância ou mudança de posição entre
audiências é registrada com o período. Frases que serviriam para qualquer participante são cortadas,
e o perfil nunca descreve o que a pessoa não disse. O comprimento é proporcional ao material, com
teto de 800 palavras; o teto é só uma instrução do prompt, e a geração não o impõe.

Duas regras de escrita decorrem de o perfil ser destinado a um codificador de texto: uma estrutura
fixa, com rótulos como ``Temas:'' e ``Posições:'', repetiria o mesmo texto em todos os perfis e
aproximaria os vetores de pessoas diferentes pela forma, e a proibição de descrever ausências evita que
frases como ``não menciona temas de saúde'' levem ao texto termos que a pessoa nunca usou. As regras
sobre o cabeçalho, sobre posições de terceiros e sobre siglas e nomes foram acrescentadas depois da
leitura de perfis gerados durante o desenvolvimento do prompt, em que o assunto da audiência aparecia
como posição do ator, a proposta de um convidado resumida pelo presidente aparecia como proposta dele,
e siglas e sobrenomes eram completados com conhecimento externo às falas.

Os dois prompts formam uma conversa de duas mensagens, formatada pelo \textit{chat template} do próprio
LLM, sem formatação adicional. A geração usa amostragem com temperatura 0,2, \textit{top-p} 0,9, até
3.000 tokens de saída e semente fixada antes de cada geração, de modo que o perfil de um ator não
depende da ordem em que os atores são processados. Uma resposta que atinge o limite de tokens sem
terminar conta como falha e é descartada. Cada perfil guarda o modelo, a versão dos prompts, as
audiências usadas e o número de tokens de entrada e de saída. O modelo gerador é <modelo>.

% Pendente: trecho de um perfil real (Quadro) depois da rodada de geração com o modelo final.

\subsection{<Objetivo final: a escrever>}
\label{sec:metodo-objetivo}

<Uso conjunto das UDVs e dos perfis. Seção pendente até o objetivo final ser definido.>
